% !TEX encoding = UTF-8
% !TEX TS-program = pdflatex
% !TEX root = ../tesi.tex

%**************************************************************
\chapter{Introduction}
\label{cap:introduction}
%**************************************************************

%**************************************************************
\section{Directional modulation demand}

\textcolor{ForestGreen}{GB: I checked all points, and most of them are solved. Only a few final remarks:\\
- Fig. 4.6: since $B_{e,\min}$ has been defined some time ago (page 24), I would add it at the end of the first sentence of Sect. 4.2.5: ``the minimum BER of all possible eavesdropper angles $B_{e,\min}$''. The same for the legitimate user BER and the secrecy BER, if you used a symbol for them.\\
- Since Figs. 4.16 and 4.10 (that you mention to justify the potentially odd behavior of Fig. 4.6) come at a later stage, I would say: ``not an artifact of incomplete convergence, as will be discussed in Section 4.3'', and remove the next sentence, from ``Figure 4.16 examines'' to ``it does not'' (you could use these sentences later one when commenting the figures, if needed);\\
- Fig. 4.19(b): I still don't see the shaded regions, I don't know if it's a problem of the pdf reader I'm using;\\
\\
Minor comments:\\
- page 43: please arrange the position of the figures so as to avoid the empty half page;\\
- caption of Fig. 4.40: ``dark ='' -> ``dark:'', ``light ='' -> ``light:'';\\
- page 71, third bullet: ``Uniform Planar Array'' should be ``Uniform planar array'';\\
- page 71, third bullet: use \texttt{Gls\{MARL\}} to get the first letter to be capital;\\
- page 71, last bullet: ``Figure'' should be ``Figures'' (you could define \texttt{figuresname} as ``Figures'');\\
- I would also add the following two references when introducing DDM (see commented text below for the .bib entries):\\
  - V. Pellegrini, F. Principe, G. de Mauro, R. Guidi, V. Martorelli, and R. Cioni, “Cryptographically secure radios based on directional modulation,” in Proc. IEEE Intl. Conf. Acoustics, Speech and Signal Process. (ICASSP), Florence, Italy, 2014.\\
  - V. Pellegrini, F. Principe, R. Guidi, G. Scozza, and G. de Mauro, “Adding PHY-layer crypto to COFDM radios through a large array with directional modulation,” in Proc. IEEE Global Commun. Conf. (GLOBECOM), Abu Dhabi, United Arab Emirates, 2018.
}
  % @inproceedings{pellegrini_2014,
  %   author    = {V. Pellegrini and F. Principe and G. {de Mauro} and R. Guidi and V. Martorelli and R. Cioni},
  %   title     = {Cryptographically secure radios based on directional modulation},
  %   booktitle = {Proc. IEEE Intl. Conf. Acoustics, Speech and Signal Process. (ICASSP)},
  %   address   = {Florence, Italy},
  %   year      = {2014}
  % }

  % @inproceedings{pellegrini_2018,
  %   author    = {V. Pellegrini and F. Principe and R. Guidi and G. Scozza and G. {de Mauro}},
  %   title     = {Adding {PHY}-Layer Crypto to {COFDM} Radios through a Large Array with Directional Modulation},
  %   booktitle = {Proc. IEEE Global Commun. Conf. (GLOBECOM)},
  %   address   = {Abu Dhabi, United Arab Emirates},
  %   year      = {2018}
  % }


Due to the broadcast nature of wireless communication channels, ensuring confidentiality and integrity is a fundamental challenge. Traditionally, the majority of security guarantees are implemented through cryptographic-based mechanisms in the upper layers of transmission. However,\gls{PLS}\cite{physec} has emerged as a complementary approach that exploits the inherent characteristics of the wireless channel between a legitimate transmitter (Alice) and the legitimate receiver (Bob) to secure the transmission.\gls{DM} is a\gls{PLS} technique that utilizes the spatial characteristics of the legitimate wireless channel such that a signal sent in a desired direction will be correctly demodulated only along a desired spatial direction while appearing disrupted in all other directions.\\
\\
There are important use cases that require\gls{DM} to function, such as:
\begin{enumerate}
	\item DM-aided UAV and star-terrestrial communications network;
	\item DM system with distributed and cooperative users;
	\item DM-aided V2X network;
	\item Relay-or-RIS-aided DM in range domain.
\end{enumerate}
Conventional (static)\gls{DM} schemes are vulnerable to advanced eavesdropping strategies, particularly when channel conditions are partially known or can be estimated.\gls{DDM} addresses this limitation by randomizing the array weight for each transmission, see \sectionname~\ref{sec:directional_modulation} for further details. This technique has desirable security properties but has minimal adoption due to its inability to satisfy both timing and quality constraints at the same time.
\section{Related works}

Directional modulation originated as a phased-array technique, first formalized with phase configurations computed directly for single-direction transmission and via a genetic algorithm for the harder multi-directional case \cite{daly2009dm}, and later demonstrated in hardware \cite{daly2010demo}. Subsequent work shifted toward closed-form and optimization-based synthesis, including vector-based \cite{ding2014vector} and orthogonal-vector, artificial-noise-informed \cite{ding2015orthogonal} approaches capable of synthesizing multi-beam transmitters that serve several legitimate directions simultaneously; see \cite{ding2016review} for a broader review of this line of work. The vector-based approach \cite{ding2014vector} is purely static, with no per-symbol randomization. The orthogonal-vector method \cite{ding2015orthogonal} does randomize per symbol, but only within the null space of the target direction: the signal-carrying component is identical at every symbol and must be re-solved whenever the target direction changes; this differs from the\gls{DDM} implementations described below \cite{cabrerahernandez2019sdr,cabrerahernandez2019vm,goriasilva2023lpddm}, where the complete physical excitation of the array varies per symbol.\\
\\
\gls{DDM} itself, varying the array configuration over time rather than fixing it once at design time, has also been explored, but predominantly through simple, non-learned randomization. Hardware implementations have realized\gls{DDM} via software-defined radio \cite{cabrerahernandez2019sdr} and vector-modulator architectures \cite{cabrerahernandez2019vm}, while \cite{goriasilva2023lpddm} shows that a low-power\gls{DDM} variant based on randomly switching the transmitting antenna element leaks exploitable information to an eavesdropper. In all three, the randomization rule is pre-computed offline and fixed during transmission, rather than generated by a learned policy that generalizes across target directions and jointly shapes per-symbol diversity for beamforming quality and secrecy.\\
\\
Having explored the growing demand for dynamic directional modulation, we also analyzed the current literature on\gls{AI}-assisted\gls{DM}. This survey \cite{chen2025survey} highlights the scarcity of research on the integration of\gls{AI} techniques into\gls{DM} systems and the need for more computationally efficient\gls{DDM} systems. Similar\gls{RL} approaches have been explored \cite{li2026rlris} successfully for\gls{DM}, however they do not explore the feasibility of\gls{RL} applied to\gls{DDM} systems. To the best of our knowledge, a\gls{RL} approach to\gls{DDM} has not been studied before: as discussed above, none of the existing\gls{DM} or\gls{DDM} works learns a policy for this specific problem, namely producing a fresh, stochastic per-symbol array excitation toward a fixed legitimate direction, learned through a reward signal rather than fixed by construction. \\
\\
The dynamic nature of\gls{DM} introduces constraints that significantly narrow the design space: configurations must be generated at the symbol rate and repeated transmissions toward the same direction must yield distinct configurations to preserve the stochastic character of the signal observed by an eavesdropper outside the protected angular region.\gls{RL}, framed as a contextual bandit, addresses these constraints simultaneously: it optimizes a reward signal, generalizes across angles, and produces stochastic outputs through its policy. A detailed justification of this design choice relative to alternative approaches is provided in \sectionname~\ref{sec:reinforcement_learning}.
\section{Thesis structure}
This thesis will examine first the issue of dynamic and static directional modulation in the context of\gls{PLS}, and explain how to possibly exploit a reward approach to train an\gls{AI} model that will provide\gls{DDM}. The concepts will be introduced gradually, requiring minimal prior knowledge on the topics of\gls{AI} and telecommunication. Here is a brief outline of the structure of the thesis:

\begin{itemize}
	\item \textbf{Chapter 2:} will introduce the main theoretical concepts of telecommunication and\gls{AI} model training;
	\item \textbf{Chapter 3:} will describe the tools utilized to train the model, the implementation details, and the parameters chosen;
	\item \textbf{Chapter 4:} will explain the methods used to validate the model, present and analyze the experimental results;
	\item \textbf{Chapter 5:} will contain the conclusions of this work, along with possible improvements  that can be introduced in the future.
\end{itemize}